\section{演習を別の視点から}\label{節：演習の操作を別の表現で見直す}
  \bookproblemref{practice-entrance}{9}{発展問9},
  \bookproblemref{practice-entrance}{4}{問4},
  \bookproblemref{practice-entrance}{3}{問3}の操作を,不変量・回転・母関数から見直す.

  \subsection{整数方程式と不変量}
  発展問9では,$u_0=3$, $u_1=3$, $u_{n+2}=3u_{n+1}-u_n$から,
  $(3,u_n,u_{n+1})$が$x^2+y^2+z^2=xyz$の解になることを示した.
  $n=0$では
  \[
    u_n^2+u_{n+1}^2+9=3u_nu_{n+1}
  \]
  が成り立つ.一般の$x,y$に対して
  \[
    y^2+(3y-x)^2-3y(3y-x)=x^2+y^2-3xy
  \]
  だから,左辺の差は更新の前後で変わらず$-9$に保たれる.
  その結果,$u_nu_{n+2}=u_{n+1}^2+9$となり,
  \[
    u_{n+2}=\frac{u_{n+1}^2+9}{u_n}
  \]
  とも書ける.正値性は$u_0=u_1=3$と,
  $u_2=6$以降の増加から確認できる.
  \sectionref*{節：不変量から線形の関係}の定数$1$を$9$へ変えた更新だった.
  元の問題は\bookterm{linear}{線形}の更新から保存を示し,本編の例は保存から線形の更新を取り出した.

  \subsection{中点・回転・面積}
  発展問4では,平行四辺形の各辺上で同じ比率$t$ $(0<t<1)$の分点を取り,
  新しい平行四辺形を作った.連続する辺ベクトルを$u,v$とする.
  頂点の更新を引き算すれば,
  \[
    u'=(1-t)u+tv,\qquad v'=-tu+(1-t)v
  \]
  である.

  \bookstep{正方形の中点更新}
  正方形で,向きを反時計回りに取った辺を複素数として$v=iu$と書けるときは,
  \[
    u'=((1-t)+it)u
  \]
  である.$t=\frac12$なら倍率は$\frac1{\sqrt2}$,回転角は$\frac\pi4$になる.
  これが問題の中点更新の辺長と向きに一致する.

  \bookstep{別解：平行四辺形を行列で見る}
  正方形という条件を外し,最初に求めた$u',v'$の式へ戻る.
  辺の組を混ぜる\bookterm{matrix}{行列}は
  \[
    M=\begin{pmatrix}1-t&t\\-t&1-t\end{pmatrix}
  \]
  となる.\bookterm{matrix}{行列}式は$\rho=(1-t)^2+t^2$である.
  $u=(u_1,u_2)$, $v=(v_1,v_2)$と書くと,面積は$|u_1v_2-u_2v_1|$である.更新を代入すると
  \BookMathLayout{\[
    u'_1v'_2-u'_2v'_1
    =\bigl((1-t)^2+t^2\bigr)(u_1v_2-u_2v_1).
  \]}{\[
    \begin{aligned}
      u'_1v'_2-u'_2v'_1\\
      {}=\bigl((1-t)^2+t^2\bigr)(u_1v_2-u_2v_1).
    \end{aligned}
  \]}
  $\rho>0$なので,絶対値を取った面積も一回ごとに$\rho$倍となる.
  この行列式の倍率は,\bookterm{eigenvalue}{固有値}$1-t\pm it$の絶対値の二乗でもある.

  \begin{bookpoint}
    \textbf{\bookterm{matrix}{行列}が作用する対象}\par
    辺の組を混ぜる\bookterm{matrix}{行列}を,平面上の相似変換と読むには条件が要る.
    面積の倍率は一般の平行四辺形に使えるが,上の回転の説明には直交・等長の辺を使った.
  \end{bookpoint}

  \subsection{一定になる畳み込み和}
  発展問3の$c_0=1$, $c_{n+1}=\frac{2n+1}{2n+2}c_n$から,
  \[
    c_n=\frac1{4^n}\binom{2n}{n}
  \]
  を得た.その\bookterm{generating}{母関数}を$C(x)=\sum_{n\ge0}c_nx^n$とする.
  $C(x)^2$の$x^n$の係数が,求める和$\sum_{k=0}^n c_kc_{n-k}$である.
  形式的な微分$x^n\mapsto nx^{n-1}$を使う.関数としての収束は仮定しない.
  元の\bookterm{recurrence}{漸化式}を$(n+1)c_{n+1}=(n+\frac12)c_n$と書けば,
  \[
    (1-x)C'(x)=\frac12C(x)
  \]
  を得る.積の係数は有限和なので,係数を比較することで通常の積の微分法則も使える.
  $D(x)=C(x)^2=\sum_{n\ge0}d_nx^n$と置くと,
  \[
    (1-x)D'(x)=2C(x)(1-x)C'(x)=D(x).
  \]
  $x^n$の係数を比較すれば$(n+1)d_{n+1}-nd_n=d_n$,すなわち$d_{n+1}=d_n$である.
  $d_0=c_0^2=1$なので,
  \[
    C(x)^2=\frac1{1-x}=\sum_{n\ge0}x^n
  \]
  の$x^n$の係数から
  \[
    \sum_{k=0}^n c_kc_{n-k}=1
  \]
  を得る.有限和の変形と母関数の積は,同じ\bookterm{convolution}{畳み込み}の二つの扱い方である.

  \begin{bookpoint}
    \textbf{別解の使い分け}\par
    別の表現は,和・面積・不変量の意味を見せる.元の条件と答えへ戻って対応を確かめる.
  \end{bookpoint}
